% !TEX root = ../main.tex
\chapteropening{%
A spacetime engineer designs geometries, so we need a precise language for
geometry. This chapter builds it from spacetime diagrams up to curvature.

With it you will be able to compute what a traveller would experience in any
spacetime: the ticking of clocks, the paths of light, the motion of falling
bodies and the stretching forces called tides.}{%
\item draw and read spacetime diagrams, worldlines and light cones;
\item compute the interval between two events and say what can influence
  what;
\item compute the time a clock records along any path;
\item use a metric in any coordinates, and tell a feature of the coordinates
  from a feature of the geometry;
\item compute clock rates in weak gravity and near a black hole, and picture a
  black hole as flowing space;
\item describe free fall as motion along a geodesic, and curvature as the
  source of tides.}
\chapter{The Geometry of Spacetime}\label{ch:geometry}

% ===========================================================================
\section{What every observer agrees on}\label{sec:interval}

\begin{thought}{Has Betelgeuse exploded?}
Betelgeuse, the red supergiant at Orion's shoulder, is roughly 550 light-years
away, and one day it will explode as a supernova. Suppose that, in the frame
in which the Earth and the star are at rest, it explodes at the very moment you
read this sentence. At that moment a ship passes the Earth, heading for
Betelgeuse at half the speed of light. Much later, working back from what they
have seen, the crew concludes that the explosion happened 318 years before
they passed you. Could they have seen it in time to warn you? And how far away
was the star, by their reckoning, when it exploded?
\end{thought}

An \emph{event} is a place at an instant: a firecracker going off at a
particular point and time. We label events with four coordinates, a time $t$
and three positions $x$, $y$, $z$. A particle traces out a continuous sequence
of events called its \emph{worldline}. A \emph{spacetime diagram} plots
worldlines with time running up the page and one direction of space across it
(\cref{fig:spacetime-diagram}). A particle at rest has a vertical worldline; a
moving particle has a tilted one, and the faster it moves the more its
worldline leans over.

\begin{figure}[tb]
\centering
\includegraphics{ch02-spacetime-diagram}
\caption{A spacetime diagram, with time running up the page. A particle at
rest has a vertical worldline and a moving particle a tilted one; a particle
that speeds up and slows down has a curved one. With $c = 1$, light travels
along lines at \ang{45}.}
\label{fig:spacetime-diagram}
\end{figure}

It pays to measure time and distance in the same units. A light-second is the
distance light travels in one second, about \SI{3e8}{\metre}. If we measure
distance in light-seconds and time in seconds, or time in metres (one metre
of time is the time light takes to travel a metre, about
\SI{3.3}{\nano\second}), then the speed of light is exactly $c = 1$ and
velocities become pure numbers: a speed of $0.5$ is half the speed of light.
From here on we use these units, and we restore the factors of $c$ whenever we
need a number in everyday units. Light worldlines then run at \ang{45} on every
spacetime diagram.

In ordinary space, two surveyors who use differently oriented axes disagree
about the coordinates of two points but agree on the distance between them:
$\Delta x^{2} + \Delta y^{2} + \Delta z^{2}$ is the same for both. Special
relativity has an exact analogue. Observers moving relative to one another
disagree about time differences and about distances, yet all of them agree on
the \emph{interval}
\begin{equation}
  \Delta s^{2} = -\Delta t^{2} + \Delta x^{2} + \Delta y^{2} + \Delta z^{2}.
  \label{eq:interval}
\end{equation}
The minus sign in front of the time difference is the whole difference between
the geometry of spacetime and the geometry of space. It sorts pairs of events
into three kinds. The interval is \textbf{timelike} if $\Delta s^{2} < 0$: a
clock can travel from one event to the other at less than the speed of light.
It is \textbf{null}, or lightlike, if $\Delta s^{2} = 0$: a light signal can
travel from one to the other. It is \textbf{spacelike} if
$\Delta s^{2} > 0$: nothing, not even light, can travel from one to the other.

The events that light can reach from an event $A$ form a cone, the
\emph{future light cone} of $A$ (\cref{fig:light-cone}). Events inside it can
be influenced by what happens at $A$; events outside it cannot. The events
from which light can reach $A$ form its past light cone. The light cones of all
events together make up the \emph{causal structure} of spacetime: they tell
you what can affect what. Every question about faster-than-light travel is, at
bottom, a question about light cones.

The interval also settles when the order of two events is shared. For a
timelike or null pair, every observer agrees which came first, since a signal
could pass from one to the other. For a spacelike pair, observers in relative
motion can disagree about the order, and none of them is wrong;
\cref{prob:simultaneity} shows how.

\begin{figure}[tb]
\centering
\includegraphics{ch02-light-cone}
\caption{The light cone of an event $A$, with one dimension of space
suppressed. Events inside the future cone can be influenced by $A$, and events
inside the past cone can influence it. Events outside both cones are
spacelike from $A$: no signal connects them to it.}
\label{fig:light-cone}
\end{figure}

\begin{checkpoint}[label=chk:classify]
Event $B$ happens \SI{3}{\second} after event $A$ and 4~light-seconds away
from it. Event $C$ happens \SI{5}{\second} after $A$, at the same place as $B$.
A probe travels from $A$ to $C$ at constant velocity: how fast does it go, and
how much time does its clock record? Could any probe be present at both $A$
and $B$? For an observer who finds $A$ and $B$ simultaneous, how far apart are
they?
\end{checkpoint}

\begin{revisited}{Has Betelgeuse exploded?}
No warning was possible. In the Earth's frame the two events, you reading and
the star exploding, are 550 light-years apart with no time between them, so
$\Delta s^{2} = (\text{550 light-years})^{2}$: a spacelike interval. The crew
must find the same interval. With $\Delta t' = 318$ years,
$\Delta x'^{2} = 550^{2} + 318^{2}$ in light-years squared, so by their
reckoning the star was 635 light-years away when it exploded, not the 476 that
length contraction gives for its distance as they pass you. Its light needs
635 years to reach them and has had only 318, so it is still some 318
light-years away when they pass you. The crew places the explosion in the
past, you place it in the present, and an observer flying away from the star
would place it in the future (\cref{prob:simultaneity}). For events with a
spacelike interval the order in time is not the same for everyone, and no
observer's version is privileged; yet no one can use the difference to send a
signal.
\end{revisited}

% ===========================================================================
\section{The time a clock keeps}\label{sec:proper-time}

\begin{thought}{The twins}
Twins say goodbye on Earth. One boards a ship to Proxima Centauri, 4.24
light-years away, travelling at 80 percent of the speed of light; she turns
round in a single day and comes straight back at the same speed. On each leg
of the trip, each twin measures the other's clock running slow, so the
situation looks symmetric. When they meet again, who is younger, and by how
much? A common answer blames the turnaround, since only the traveller
accelerates. If that were right, how would the answer change for a star ten
times as far away, with the same one-day turnaround?
\end{thought}

A clock carried along a worldline measures the \emph{proper time} $\tau$ along
it. For a clock at rest, proper time is just coordinate time. For a moving
clock the interval does the bookkeeping. Along a short piece of a timelike
worldline the interval is negative, and the clock records
\begin{equation}
  \dd\tau^{2} = -\dd s^{2} = \dd t^{2} - \dd x^{2} - \dd y^{2} - \dd z^{2}.
  \label{eq:proper-time}
\end{equation}
Dividing by $\dd t^{2}$ gives $\dd\tau/\dd t = \sqrt{1 - v^{2}}$, where $v$ is
the clock's speed, and the proper time along a whole worldline is the sum of
the pieces:
\begin{equation}
  \tau = \int \sqrt{1 - v(t)^{2}}\;\dd t .
  \label{eq:proper-time-integral}
\end{equation}
This is the time dilation of special relativity written as a geometric length:
proper time is to a worldline what length is to a curve in space.

\begin{example}[label=ex:muon]{Muons from the sky}
Cosmic rays striking the upper atmosphere create muons about
\SI{15}{\kilo\metre} up, and many of them head for the ground at
$v = 0.9995$. A muon at rest decays after \SI{2.2}{\micro\second} on average.
How much time passes on a muon's own clock during the trip down?

\solution In the Earth's frame the trip takes
\begin{equation*}
  \Delta t = \frac{\SI{15}{\kilo\metre}}{0.9995\,c} = \SI{50.1}{\micro\second}.
\end{equation*}
The muon moves at constant speed, so \cref{eq:proper-time-integral} gives
\begin{equation*}
  \Delta\tau = \sqrt{1 - 0.9995^{2}}\;\Delta t = 0.0316 \times \SI{50.1}{\micro\second}
  = \SI{1.58}{\micro\second}.
\end{equation*}

\discussion The trip takes less than one mean lifetime on the muon's clock, so
about half the muons survive it. Measured in Earth time the trip lasts 23
lifetimes, and without time dilation only about one muon in eight billion
would reach the ground. Detectors on the ground see the larger number.
\end{example}

\begin{keyidea}
Between two events, the straight worldline of a freely moving clock records
the longest proper time. Every other worldline between the same two events,
however it bends, records less.
\end{keyidea}

\begin{checkpoint}[label=chk:zigzag]
Two events happen at the same place, ten years apart. Describe a worldline
between them along which a clock records as little time as you like. Is there
a worldline of \emph{least} proper time?
\end{checkpoint}

\begin{revisited}{The twins}
The traveller is younger, by 4.24 years. In the Earth's frame the round trip
covers 8.48 light-years at speed 0.8 and takes 10.6 years; on her clock it
takes $0.6 \times 10.6 = 6.36$ years. For a star ten times as far away the gap
grows tenfold, to 42.4 years, although the turnaround is unchanged: the
ageing builds up along the whole trip, not at the corner. Between the same
two events, the traveller's bent worldline is shorter than the Earth twin's
straight one (\cref{fig:twins}). The symmetry is only apparent, and the twins
can check it by sending each other a signal every year. Each receives the
other's signals at a third of the rate they are sent while the two move
apart, and at three times the rate while they approach. The traveller
switches from one to the other at her turnaround, halfway through her trip,
and receives $1.06 + 9.54 = 10.6$ years' worth of Earth's signals. The Earth
twin receives slow signals until the light from the turnaround arrives, 9.54
years after the start, and fast ones only for the last 1.06 years:
$3.18 + 3.18 = 6.36$ years' worth.
\end{revisited}

\begin{figure}[tb]
\centering
\includegraphics{ch02-twins}
\caption{The twins' worldlines. Each tick marks one year on that twin's own
clock. The straight worldline of the twin who stays home collects 10.6 ticks;
the bent worldline of the traveller collects 6.36.}
\label{fig:twins}
\end{figure}

% ===========================================================================
\section{The metric: measuring with any grid}\label{sec:metric}

\begin{thought}{A current faster than light}
In \cref{ch:what-is}, a radio command sent forward inside a warp bubble
stalled in the bubble's wall, where space streamed backward past the ship at
the speed of light. Suppose instead that space streams past a grid of
coordinates at twice the speed of light, at the same rate everywhere. Two
probes drift with the stream, one light-second apart, and the downstream probe
radios the upstream one. Does the message ever arrive?
\end{thought}

\Cref{eq:interval} holds for inertial observers using Cartesian coordinates. We
need a version that works in any coordinates, and later in curved spacetime.
The idea is to give the interval between \emph{neighbouring} events as a
formula in whatever coordinates we are using.

We label the four coordinates with an index that runs from 0 to 3:
$x^{0} = t$, $x^{1} = x$, $x^{2} = y$, $x^{3} = z$. The raised index is a
label, not a power, and Greek letters $\mu, \nu, \dots$ stand for any of the
four values. The interval between neighbouring events is then
\begin{equation}
  \dd s^{2} = g_{\mu\nu}\,\dd x^{\mu}\,\dd x^{\nu},
  \label{eq:metric}
\end{equation}
where we use the \emph{summation convention}: an index that appears once up and
once down in a term is summed from 0 to 3, so the right side stands for
$\sum_{\mu}\sum_{\nu} g_{\mu\nu}\,\dd x^{\mu}\dd x^{\nu}$. The sixteen numbers
$g_{\mu\nu}$ form a symmetric $4\times4$ matrix called the \emph{metric}. For
flat spacetime in Cartesian coordinates, comparison with \cref{eq:interval}
gives the Minkowski metric
\begin{equation*}
  g_{\mu\nu} = \eta_{\mu\nu} =
  \begin{pmatrix} -1 & 0 & 0 & 0\\ 0 & 1 & 0 & 0\\ 0 & 0 & 1 & 0\\ 0 & 0 & 0 & 1\end{pmatrix}.
\end{equation*}

The same flat spacetime looks different in other coordinates. In spherical
coordinates the interval reads
\begin{equation}
  \dd s^{2} = -\dd t^{2} + \dd r^{2} + r^{2}\,\dd\Omega^{2},
  \label{eq:flat-spherical}
\end{equation}
with $\dd\Omega^{2} = \dd\theta^{2} + \sin^{2}\theta\,\dd\phi^{2}$. The metric
components now depend on position, yet nothing physical has changed.
Coordinates are labels we attach to events, and much of the skill of working
with metrics lies in telling which features of a metric belong to the labels
and which belong to the geometry. The next example shows how dramatic a pure
effect of the labels can look.

\begin{example}[label=ex:flow]{Flowing coordinates}
Consider the metric
\begin{equation}
  \dd s^{2} = -\dd t^{2} + \bigl(\dd x - v_{0}\,\dd t\bigr)^{2} + \dd y^{2} + \dd z^{2},
  \label{eq:flow-metric}
\end{equation}
with $v_{0}$ a constant. Find the coordinate speed $\dd x/\dd t$ of light moving
along $x$, and show that this metric describes flat spacetime.

\solution Light has $\dd s^{2} = 0$. Along $x$ we set $\dd y = \dd z = 0$ and
find $(\dd x - v_{0}\,\dd t)^{2} = \dd t^{2}$, so $\dd x - v_{0}\,\dd t = \pm\dd t$
and
\begin{equation*}
  \frac{\dd x}{\dd t} = v_{0} \pm 1.
\end{equation*}
Now change coordinates to $x' = x - v_{0}t$, keeping $t$, $y$ and $z$. Then
$\dd x' = \dd x - v_{0}\,\dd t$, and the metric becomes
$\dd s^{2} = -\dd t^{2} + \dd x'^{2} + \dd y^{2} + \dd z^{2}$, the Minkowski
metric.

\discussion This is flat spacetime described with a grid whose points slide
along $-x$ at speed $v_{0}$ past inertial observers. The warp drives of
\cref{ch:what-is} have metrics of exactly this form, with $v_{0}$ replaced by a
function that varies from place to place, and then the effect is no longer
only a matter of labels.
\end{example}

The form \cref{eq:flow-metric} gives us a picture that runs through this book:
space \emph{flowing} past a grid of coordinates. Observers carried along by the
flow feel no force. The flow speed can take any value. Where it varies from
place to place the geometry may be genuinely curved, though not always: a flow
that rotates rigidly, for instance, is still flat spacetime in rotating
coordinates, and only a calculation of curvature settles the question.

The tangent to a worldline, measured per unit proper time, is the
\emph{four-velocity} $u^{\mu} = \dd x^{\mu}/\dd\tau$. Dividing \cref{eq:metric}
by $\dd\tau^{2} = -\dd s^{2}$ gives its normalization,
\begin{equation}
  g_{\mu\nu}\,u^{\mu}u^{\nu} = -1.
  \label{eq:four-velocity}
\end{equation}
For a clock at rest in flat spacetime, $u^{\mu} = (1,0,0,0)$; for a clock moving
at speed $v$ along $x$, $u^{\mu} = \gamma\,(1, v, 0, 0)$ with
$\gamma = 1/\sqrt{1-v^{2}}$. In \cref{ch:matter-geometry} four-velocities tell
us what an observer measures.

\begin{checkpoint}[label=chk:grid-clock]
In the flowing coordinates of \cref{ex:flow} with $v_{0} = 0.6$, a clock rides
with one grid point, at fixed $x$. How much time does it record while the
coordinate time $t$ advances by ten years? Compare with the time dilation of a
clock moving at speed $0.6$.
\end{checkpoint}

\begin{revisited}{A current faster than light}
It arrives after one second, exactly as if there were no stream at all. A
stream that is the same everywhere is flat spacetime described with sliding
coordinates: it is \cref{eq:flow-metric} with $v_{0} = 2$, and in the
coordinates $x' = x - v_{0}t$ of \cref{ex:flow} the metric is Minkowski and
both probes sit at rest. In the grid coordinates the numbers look alarming.
The probes move at $\dd x/\dd t = 2$, and the message, sent upstream, moves at
$v_{0} - 1 = +1$, carried downstream by the stream; yet the upstream probe
gains on it and meets it one second later. What stalled the command in the
warp bubble is something no change of coordinates removes: a stream that
changes, across the wall, from nothing to faster than light. A horizon needs
flow that reaches the speed of light next to slower flow, where an observer
can hold position against it, as in the river of space two topics ahead.
\end{revisited}

% ===========================================================================
\section{Clocks in a gravitational field}\label{sec:clocks}

\begin{thought}{Light climbing a tower}
\Cref{ch:what-is} told you that a clock at the top of a tower runs fast
compared with one at its foot, by the fraction $gh/c^{2}$. A laser at the foot
shines a steady beam up to a detector at the top. Each photon carries the
energy $E = hf$, where $f$ is the light's frequency and $h$ is Planck's
constant, and energy has weight, so each photon loses the fraction $gh/c^{2}$
of its energy on the way up. Someone reasons that the detector must find the
frequency lower by $2gh/c^{2}$ in all: once because the light arrives with
less energy, and again because the detector's clock runs fast. Is that right?
\end{thought}

In curved spacetime no choice of coordinates puts the metric into the
Minkowski form everywhere at once. The metric still gives the interval between
neighbouring events through \cref{eq:metric}, and the proper time along a
worldline is still $\tau = \int\sqrt{-\dd s^{2}}$. Everything we did above
carries over; only the metric changes.

Around a planet or a star the gravitational field is weak and the metric is
very close to Minkowski. To an excellent approximation,
\begin{equation}
  \begin{split}
  \dd s^{2} = {}&-(1 + 2\Phi)\,\dd t^{2}\\
  &+ (1 - 2\Phi)\bigl(\dd x^{2} + \dd y^{2} + \dd z^{2}\bigr),
  \end{split}
  \label{eq:weak-field}
\end{equation}
where $\Phi$ is the Newtonian gravitational potential divided by $c^{2}$. For a
spherical mass $M$, $\Phi = -GM/(c^{2}r)$, small and negative:
$\Phi \approx -7\times10^{-10}$ at the Earth's surface, as in
\cref{ex:earth-dent}. The approximation holds wherever $|\Phi| \ll 1$.

A clock moving slowly at speed $v$ through this spacetime records, from
\cref{eq:weak-field},
\begin{equation}
  \begin{split}
  \frac{\dd\tau}{\dd t} &= \sqrt{1 + 2\Phi - (1 - 2\Phi)\,v^{2}}\\
  &\approx 1 + \Phi - \tfrac{1}{2}v^{2},
  \end{split}
  \label{eq:clock-rate}
\end{equation}
keeping only the leading terms in the small quantities $\Phi$ and $v^{2}$. The
two corrections have a clear meaning: a clock deeper in the potential (more
negative $\Phi$) runs slower, and a moving clock runs slower.

The same clock rates set what a receiver finds in a signal. The metric
\cref{eq:weak-field} does not change with $t$, so every wave crest sent from
one place to another takes the same coordinate time for the trip: crests that
leave a transmitter a coordinate interval $\Delta t$ apart arrive $\Delta t$
apart. Each clock converts $\Delta t$ into its own time at its own rate. A
receiver higher up, where clocks run faster, counts more of its own time
between crests than the transmitter did, and so measures a lower frequency,
lower by the fractional difference in potential, $\Delta\Phi$. This is the
\emph{gravitational redshift}. Near the ground $\Delta\Phi = gh/c^{2}$ for a
height difference $h$. In 1960 Robert Pound and Glen Rebka measured it with
gamma rays sent up and down a tower \SI{22.6}{\metre} tall at Harvard, where the
shift is only $2.5\times10^{-15}$, and their result agreed with the prediction
to within ten percent. Today's optical clocks detect the change in rate when
one of them is raised by just \SI{33}{\centi\metre}, a shift of four parts in
$10^{17}$.

\begin{example}[label=ex:gps, unbreakable]{Clocks on the ground and in orbit}
Navigation satellites orbit at a radius
$r_{\mathrm{sat}} = \SI{26560}{\kilo\metre}$. Take
$GM = \SI{3.986e14}{\metre\cubed\per\second\squared}$ and
$R = \SI{6371}{\kilo\metre}$ for the Earth, and ignore its rotation. How much
does a satellite clock gain each day over a clock on the ground?

\solution The ground clock is at rest at $r = R$, so its rate is $1 + \Phi(R)$.
The satellite clock moves on a circular orbit at speed
$v = \sqrt{GM/r_{\mathrm{sat}}} = \SI{3.87}{\kilo\metre\per\second}$. In SI
units the difference in rates is
\begin{equation*}
  \frac{GM}{c^{2}}\left(\frac{1}{R} - \frac{1}{r_{\mathrm{sat}}}\right)
  - \frac{v^{2}}{2c^{2}},
\end{equation*}
a gain from height of $5.29\times10^{-10}$ and a loss from speed of
$8.35\times10^{-11}$. Over one day, \SI{86400}{\second}, the satellite clock
gains \SI{45.7}{\micro\second} from its height and loses
\SI{7.2}{\micro\second} from its speed, a net gain of
\SI{38.5}{\micro\second}.

\discussion These are the numbers behind \cref{fig:clock-offsets}. Repeating
the calculation for every orbital radius draws the three curves of that
figure (\cref{prob:gps-curves}).
\end{example}

For a non-rotating mass $M$, and with no weak-field approximation at all, the
geometry outside the mass is the \emph{Schwarzschild metric},
\begin{equation}
  \begin{split}
  \dd s^{2} = {}&-\Bigl(1 - \frac{r_{\mathrm{s}}}{r}\Bigr)\dd t^{2}
  + \frac{\dd r^{2}}{1 - r_{\mathrm{s}}/r}\\
  &+ r^{2}\,\dd\Omega^{2},
  \end{split}
  \label{eq:schwarzschild}
\end{equation}
with $r_{\mathrm{s}} = 2GM/c^{2}$. Here $r$ is defined so that a sphere at
radius $r$ has area $4\pi r^{2}$. A clock held at fixed $r$ records
$\dd\tau/\dd t = \sqrt{1 - r_{\mathrm{s}}/r}$ relative to a clock far away. For
the Sun $r_{\mathrm{s}} \approx \SI{2.95}{\kilo\metre}$, so the effect at its
surface is tiny, but a clock on the surface of a neutron star runs at about
0.81 of the rate of a distant clock (\cref{prob:neutron-star}).

At $r = r_{\mathrm{s}}$ the clock rate would fall to zero: no clock can be held
at rest there. The sphere $r = r_{\mathrm{s}}$ is the \emph{event horizon} of a
black hole. The metric \cref{eq:schwarzschild} also misbehaves there, where
$g_{rr}$ becomes infinite, but that part is an artefact of the labels, as the
next topic shows.

\begin{checkpoint}[label=chk:earth-centre]
A clock sits in a small cavity at the centre of the Earth. It feels no pull at
all, just like a clock far out in space. Does it keep the same rate as the
distant clock? Is it faster or slower than a clock on the Earth's surface?
\end{checkpoint}

\begin{revisited}{Light climbing a tower}
No: the frequency is lower by $gh/c^{2}$, once. The geometry does not change
with time, so crests reach the top exactly as often as they leave the bottom,
counted in the coordinate time $t$. The only difference between the two ends
is the rate of their clocks. The top clock runs faster, so it counts a longer
time between crests and reports a frequency lower by $gh/c^{2}$, and with
$E = hf$ that lower frequency \emph{is} the lost energy. The energy loss and
the fast clock are one effect described twice, and adding them counts it
twice. Pound and Rebka measured $1.05 \pm 0.10$ times the prediction, which
rules out the doubled value.
\end{revisited}

% ===========================================================================
\section{Space as a flowing river}\label{sec:river}

\begin{thought}{Standing still in the river}
A probe hangs on a long cable at twice the horizon radius of a black hole,
$r = 2r_{\mathrm{s}}$. Its clock runs at $\sqrt{1 - r_{\mathrm{s}}/r} = 0.71$ of
the rate of clocks far away. A second probe, dropped from rest far away, falls
past it at 71~percent of the speed of light. Moving clocks run slow; but which
of these two probes is moving? Is it a coincidence that the two numbers match?
\end{thought}

Introduce a new time coordinate $T$: the time on the clocks of observers who
fall inward from rest far away. In these coordinates the Schwarzschild geometry
takes the Painlevé--Gullstrand form,
\begin{equation}
  \dd s^{2} = -\dd T^{2} + \Bigl(\dd r + \sqrt{\tfrac{r_{\mathrm{s}}}{r}}\;\dd T\Bigr)^{2}
  + r^{2}\,\dd\Omega^{2}.
  \label{eq:river}
\end{equation}
Compare it with \cref{eq:flow-metric}. It has the same form, with space flowing
\emph{inward} at the speed $\sqrt{r_{\mathrm{s}}/r}$, which grows as $r$ shrinks
(\cref{fig:river}). Nothing is singular at $r = r_{\mathrm{s}}$ any more. Radial
light obeys $\dd s^{2} = 0$, which gives
\begin{equation}
  \frac{\dd r}{\dd T} = -\sqrt{\frac{r_{\mathrm{s}}}{r}} \pm 1 :
  \label{eq:river-light}
\end{equation}
light always moves at speed 1 relative to the space around it, and the space
carries it inward. Like a fish in a river that runs ever faster toward a
waterfall, outgoing light makes headway only where the flow is slower than its
own speed. At $r = r_{\mathrm{s}}$ the flow reaches the speed of light, and
inside that sphere even outgoing light is swept inward: the sphere is the
horizon. This picture, developed by Andrew Hamilton and Jason Lisle as the
\emph{river model} of black holes, will be our main intuition for geometries
whose space flows.

\begin{figure}[tb]
\centering
\includegraphics{ch02-river}
\caption{The Schwarzschild geometry pictured as space flowing inward
(\cref{eq:river}). Top: the inward flow speed $\sqrt{r_{\mathrm{s}}/r}$
(blue) reaches the speed of light (dashed) at the horizon. Bottom: the flow
drawn to scale (blue), and the net speed of outgoing light (orange), which
falls to zero at the horizon and turns inward inside it.}
\label{fig:river}
\end{figure}

\begin{example}[label=ex:earth-river]{The river at the Earth's surface}
In SI units the flow speed is $v = c\sqrt{r_{\mathrm{s}}/r} = \sqrt{2GM/r}$. How
fast does space flow inward at the surface of the Earth?

\solution With $GM = \SI{3.986e14}{\metre\cubed\per\second\squared}$ and
$r = \SI{6.371e6}{\metre}$,
\begin{equation*}
  v = \sqrt{\frac{2 \times 3.986\times10^{14}}{6.371\times10^{6}}}\ \si{\metre\per\second}
  = \SI{11.2}{\kilo\metre\per\second}.
\end{equation*}

\discussion This is the Earth's escape speed. The river of space flows at the
Newtonian escape speed everywhere, which is why an observer falling from rest
far away, carried along by the river, arrives at each radius with exactly that
speed.
\end{example}

\begin{history}{dark stars}
In a paper read to the Royal Society in 1783, John Michell asked what gravity
does to light leaving a star. Treating light as particles obeying Newton's
laws, he found that a star with the Sun's density and 500 times its radius
would have an escape speed greater than the speed of light, so its light
could never leave it. Such \enquote{dark stars} could be found only through
their gravity on companions. The radius at which Newton's escape speed reaches
$c$ is $2GM/c^{2}$, the same radius at which general relativity places the
horizon.
\end{history}

\begin{revisited}{Standing still in the river}
It is no coincidence: in the river picture the hanging probe is the one that
moves. Space flows inward past it at $\sqrt{r_{\mathrm{s}}/r}$, which is 0.71 at
$r = 2r_{\mathrm{s}}$, while the falling probe drifts with the flow, at rest
relative to the space around it, and its clock keeps pace with clocks far
away. By \cref{eq:river}, a clock held at fixed $r$, with $\dd r = 0$, records
$\dd\tau = \sqrt{1 - r_{\mathrm{s}}/r}\,\dd T$: exactly the time dilation of a
clock moving at $\sqrt{r_{\mathrm{s}}/r}$ through the local space. The
gravitational slowing of a hanging clock is its motion upstream through
flowing space. At the Earth's surface, where the river flows at the escape
speed (\cref{ex:earth-river}), it is the 60 microseconds a day by which
ground clocks lag clocks far away. At the horizon the flow reaches the speed
of light, and no probe can hang there at all.
\end{revisited}

% ===========================================================================
\section{Falling is straight}\label{sec:geodesics}

\begin{thought}{The thrown ball}
You throw a ball straight up and catch it two seconds later. On its way the
ball rises to where clocks run faster, but it also moves relative to you, and
moving clocks run slow. Which records more time between the throw and the
catch: a tiny clock inside the ball, or your wristwatch? Could you have thrown
the ball along some other path, back into your hand at the same moment, so
that its clock recorded even more?
\end{thought}

The key idea of \topicref{sec:proper-time} becomes, in general relativity, the
law of motion for anything on which no force other than gravity acts.

\begin{keyidea}[free fall]
A freely falling body follows a \emph{geodesic} of spacetime: a worldline that
makes the proper time stationary, and that, over any short enough stretch,
records more time than every nearby worldline between the same events.
Gravity is not a force that deflects it; it is the curvature that bends what
\enquote{straight} means.
\end{keyidea}

Making the proper time stationary leads, after a calculation you can do in
\cref{prob:geodesic}, to the \emph{geodesic equation}
\begin{equation}
  \frac{\dd^{2}x^{\mu}}{\dd\tau^{2}}
  + \Gamma^{\mu}{}_{\alpha\beta}\,\frac{\dd x^{\alpha}}{\dd\tau}\frac{\dd x^{\beta}}{\dd\tau} = 0,
  \label{eq:geodesic}
\end{equation}
where the \emph{Christoffel symbols} are built from first derivatives of the
metric:
\begin{equation}
  \Gamma^{\mu}{}_{\alpha\beta} = \tfrac{1}{2}\,g^{\mu\nu}
  \bigl(\partial_{\alpha}g_{\nu\beta} + \partial_{\beta}g_{\nu\alpha}
  - \partial_{\nu}g_{\alpha\beta}\bigr).
  \label{eq:christoffel}
\end{equation}
Here $\partial_{\alpha}$ means $\partial/\partial x^{\alpha}$, and $g^{\mu\nu}$ is
the inverse of the matrix $g_{\mu\nu}$. In flat spacetime with Cartesian
coordinates every $\Gamma$ vanishes, and \cref{eq:geodesic} says that free
bodies move in straight lines at constant speed.

\begin{example}[label=ex:newton, unbreakable]{Newton's law of gravity from geometry}
Show that for a slowly moving body in the weak-field metric
\cref{eq:weak-field}, with a potential that does not change in time, the
geodesic equation reduces to Newton's law $\dd^{2}x^{i}/\dd t^{2} = -\partial_{i}\Phi$.

\solution A slow body has $\dd x^{i}/\dd\tau \ll \dd t/\dd\tau$, so in
\cref{eq:geodesic} the term with $\alpha = \beta = 0$ dominates:
\begin{equation*}
  \frac{\dd^{2}x^{i}}{\dd\tau^{2}} \approx -\Gamma^{i}{}_{00}\left(\frac{\dd t}{\dd\tau}\right)^{2}.
\end{equation*}
With $g_{00} = -(1+2\Phi)$ independent of time, \cref{eq:christoffel} gives
$\Gamma^{i}{}_{00} = -\tfrac{1}{2}g^{ii}\partial_{i}g_{00} \approx \partial_{i}\Phi$
to first order. Since $\dd t/\dd\tau \approx 1$ for a slow body,
$\dd^{2}x^{i}/\dd t^{2} \approx -\partial_{i}\Phi$.

\discussion Newton's law of gravity comes from the $tt$ part of the metric
alone, that is, from the way the rate of clocks changes from place to place.
Objects fall toward regions where clocks run slower. This is one of the most
important intuitions in the book: the rate of clocks, and how it varies across
space, is itself a gravitational field.
\end{example}

\begin{revisited}{The thrown ball}
The ball's clock records more. Between the throw and the catch the ball falls
freely, so its worldline is a geodesic; your hand, held up by your arm,
follows a different worldline between the same two events. In the nearly
uniform field near the ground, the geodesic records more time than every other
worldline joining the two events. By
\cref{eq:clock-rate} the ball's clock gains $\int (gh - \tfrac{1}{2}v^{2})\,\dd t/c^{2}$
over your watch, and for the actual path this comes to $v_{0}^{3}/(3gc^{2})$,
about \SI{3.6e-16}{\second} for a two-second throw
(\cref{prob:thrown-ball}). A higher path gains more from height but loses more
from speed, a lower path the reverse; every other path back into your hand at
the same moment records less. Newton's law of motion is exactly the rule that
picks out the path of greatest proper time.
\end{revisited}

% ===========================================================================
\section{Tides: the signature of curvature}\label{sec:tides}

\begin{thought}{The windowless capsule}
Astronauts in a sealed, windowless capsule float weightless. They might be
drifting in deep space, far from any star, or orbiting the Earth, which is
also free fall. Is there any experiment they can do with loose objects inside
the capsule to tell which?
\end{thought}

A single freely falling observer feels nothing: in a small enough laboratory,
free fall cannot be told apart from floating in empty space. This is the
\emph{equivalence principle}, and it is the physical reason gravity can be
described as geometry at all. Its mathematical counterpart holds for any
metric: at any one event, coordinates can be chosen that make the metric look
like Minkowski's. Gravity reveals itself when you compare \emph{nearby}
free-fall paths.

Drop two balls side by side above the Earth. Both fall toward the Earth's
centre, so their paths converge and the balls drift together. Drop two balls
one above the other, and the lower one, closer to the Earth, accelerates
slightly more, so they drift apart (\cref{fig:tides}). These relative
accelerations are \emph{tides}, and no choice of coordinates removes them.
They are the true signature of curvature.

\begin{figure}[tb]
\centering
\includegraphics{ch02-tides}
\caption{Tides. A ring of freely falling particles above the Earth is
stretched along the direction to the Earth's centre and squeezed across it.
The arrows show the accelerations of the particles relative to the centre of
the ring.}
\label{fig:tides}
\end{figure}

In general relativity the relative acceleration of two nearby free-fall paths
separated by a small displacement $\xi^{\mu}$ is
\begin{equation}
  \frac{\dd^{2}\xi^{\mu}}{\dd\tau^{2}} = -R^{\mu}{}_{\alpha\nu\beta}\,u^{\alpha}\xi^{\nu}u^{\beta},
  \label{eq:geodesic-deviation}
\end{equation}
where $u^{\mu}$ is the four-velocity of the falling bodies and
$R^{\mu}{}_{\alpha\nu\beta}$ is the \emph{Riemann curvature tensor},
\begin{equation}
  \begin{split}
  R^{\mu}{}_{\alpha\nu\beta} = {}&\partial_{\nu}\Gamma^{\mu}{}_{\alpha\beta}
  - \partial_{\beta}\Gamma^{\mu}{}_{\alpha\nu}\\
  &+ \Gamma^{\mu}{}_{\nu\lambda}\Gamma^{\lambda}{}_{\alpha\beta}
  - \Gamma^{\mu}{}_{\beta\lambda}\Gamma^{\lambda}{}_{\alpha\nu}.
  \end{split}
  \label{eq:riemann}
\end{equation}
The Riemann tensor vanishes everywhere if and only if spacetime is flat. When
\cref{ex:flow} asked whether a metric describes flat spacetime, computing its
Riemann tensor would have answered directly, with no need to guess the right
change of coordinates.

In the weak-field limit, the tidal part of the Riemann tensor becomes the
familiar Newtonian tidal tensor, $R^{i}{}_{0j0} \approx \partial_{i}\partial_{j}\Phi$.
Near the Earth the stretching acceleration across a person \SI{2}{\metre} tall
is about $2GML/r^{3} \approx \SI{6e-6}{\metre\per\second\squared}$, far too
small to notice. Near a black hole of a few solar masses it would tear a person
apart. Tides are one of the first things a spacetime engineer checks on behalf
of a traveller.

\begin{checkpoint}[label=chk:tidal-volume]
A small ball of freely falling particles above the Earth is stretched along
the direction to the Earth's centre by a relative acceleration of $2GM/r^{3}$
per unit length, and squeezed along each of the two directions across it by
$GM/r^{3}$ per unit length. Does the volume of the ball begin to grow, to
shrink, or neither?
\end{checkpoint}

Two combinations of the Riemann tensor carry the information Einstein's
equation needs. The \emph{Ricci tensor} and the \emph{Ricci scalar} are
\begin{equation*}
  R_{\alpha\beta} = R^{\mu}{}_{\alpha\mu\beta}, \qquad R = g^{\alpha\beta}R_{\alpha\beta},
\end{equation*}
and the \emph{Einstein tensor} is
\begin{equation}
  G_{\alpha\beta} = R_{\alpha\beta} - \tfrac{1}{2}R\,g_{\alpha\beta}.
  \label{eq:einstein-tensor}
\end{equation}
The Einstein tensor is the \enquote{curvature of spacetime} on the left side of
Einstein's equation. \Cref{ch:matter-geometry} explains why this particular
combination appears, and uses it to compute what a geometry demands.

\begin{deeper}{tensors and coordinates}
The objects in this chapter are \emph{tensors}: quantities whose components
change in a definite way under a change of coordinates, so that an equation
between tensors holds in every coordinate system if it holds in one. The metric
is a symmetric tensor of type $(0,2)$ and the Riemann tensor has type $(1,3)$.
The Christoffel symbols are \emph{not} the components of a tensor, which is why
they can vanish at a point in one set of coordinates and not in another.
Spacetime is modelled as a four-dimensional Lorentzian manifold, and the sign
conventions here are those of Misner, Thorne and Wheeler; Wald develops the
same material in coordinate-free form.
\end{deeper}

\begin{revisited}{The windowless capsule}
Yes: they can look for tides. Near the Earth, two balls released a few metres
apart, one above the other, drift apart, with a relative acceleration of a few
millionths of a metre per second squared for every metre of separation. In a
minute, two balls released ten metres apart drift apart by a few
centimetres. In deep space the balls stay where they were
put. The relative acceleration is curvature, and no choice of frame removes
it. The equivalence principle holds only in a laboratory small enough, and an
experiment short enough, for tides to go unnoticed.
\end{revisited}

% ===========================================================================
\begin{chaptersummary}
\item With $c = 1$, time and distance share units and light moves at \ang{45}
  on a spacetime diagram.
\item The interval $\Delta s^{2} = -\Delta t^{2} + \Delta x^{2} + \Delta y^{2} + \Delta z^{2}$
  is the same for all inertial observers. It is timelike, null or spacelike,
  and light cones divide spacetime into what an event can and cannot
  influence.
\item A clock measures proper time $\tau = \int\sqrt{-\dd s^{2}}$ along its
  worldline. In flat spacetime the straight worldline between two events
  records the most proper time; in curved spacetime a free-fall worldline
  records more than every nearby worldline over any short enough stretch.
\item The metric $g_{\mu\nu}$ gives the interval between neighbouring events in
  any coordinates. Effects of the coordinates can be dramatic, as in flowing
  coordinates, where light's coordinate speed differs from 1.
\item In a weak field a clock runs at $\dd\tau/\dd t \approx 1 + \Phi - v^{2}/2$.
  Light climbing out of a potential arrives redshifted, because clocks at
  different heights tick at different rates. Near a black hole the rate of a
  clock held at rest falls to zero at the horizon.
\item The Schwarzschild geometry can be pictured as space flowing inward at the
  Newtonian escape speed; the horizon is where the flow reaches the speed of
  light. A clock held at rest is slowed by its motion upstream through the
  flow.
\item Free bodies follow geodesics, and the variation of clock rates from place
  to place acts as a gravitational field. The Riemann tensor measures tides,
  the true signature of curvature; the Ricci and Einstein tensors are built
  from it.
\end{chaptersummary}

\begin{checkanswers}
\item[\ref{chk:classify}] The probe covers 4 light-seconds in 5 seconds, at
  speed 0.8. Its clock records $\sqrt{5^{2} - 4^{2}} = 3$ seconds, since the
  interval $AC$ is timelike, $\Delta s^{2} = -25 + 16 = -9$. The interval $AB$
  is spacelike, $\Delta s^{2} = -9 + 16 = +7$, so no probe can be present at
  both. An observer who finds $A$ and $B$ simultaneous has $\Delta t' = 0$, and
  the invariance of the interval gives $\Delta x'^{2} = 7$: the events are
  $\sqrt{7} \approx 2.65$ light-seconds apart for that observer.
\item[\ref{chk:zigzag}] Zigzag back and forth at nearly the speed of light. On
  each leg $\dd\tau = \sqrt{1 - v^{2}}\,\dd t$, which approaches zero as $v \to 1$,
  so the total can be made as small as you like: at $v = 0.99$ the clock
  records 1.4 of the ten years. The limit, zero, belongs to light, which
  carries no clock. There is no worldline of least proper time, only one of
  greatest: the straight one, which records all ten years.
\item[\ref{chk:grid-clock}] With $\dd x = 0$,
  $\dd s^{2} = -(1 - v_{0}^{2})\,\dd t^{2} = -0.64\,\dd t^{2}$, so
  $\dd\tau = 0.8\,\dd t$: eight years. The grid point moves at speed 0.6 past
  the inertial observers whose clocks keep the time $t$, and
  $10\sqrt{1 - 0.6^{2}} = 8$ is exactly the time dilation of a clock moving at
  that speed. The metric has done the calculation of special relativity for
  us.
\item[\ref{chk:earth-centre}] Neither: the rate follows the potential $\Phi$,
  not the pull. Going down from the surface, the pull still points toward the
  centre, so $\Phi$ keeps falling all the way down. The centre is the bottom of
  the potential, and its clock is the slowest of the three. For an Earth of
  uniform density, $\Phi$ at the centre is 1.5 times its value at the surface,
  so the centre clock runs slow compared with a surface clock by
  $0.5\,GM/(c^{2}R) \approx 3.5\times10^{-10}$, about 11 milliseconds a year.
  The real Earth, denser toward its centre, gives somewhat more.
\item[\ref{chk:tidal-volume}] Neither. The volume's acceleration, as a
  fraction of the volume, is the sum of the relative accelerations per unit
  length along three perpendicular directions:
  $2GM/r^{3} - GM/r^{3} - GM/r^{3} = 0$. In empty space, tides change the shape
  of a ball of falling particles but, at first, not its volume. The sum is
  $-\nabla^{2}\Phi$, which vanishes wherever there is no matter;
  \cref{ch:matter-geometry} shows what matter does to it.
\end{checkanswers}

\begin{problems}
\item \diff{1} Two events are separated by $\Delta t = \SI{2}{\second}$ and
  $\Delta x = 1.5$~light-seconds. Find the interval and classify it. What proper
  time does a clock record if it travels between them at constant velocity?
\item \diff{2}\label{prob:simultaneity} \textbf{Whose now?} The crew of the ship
  in the Betelgeuse thought experiment flies past the Earth toward the star at
  speed 0.5. Using the Lorentz transformation $t' = \gamma(t - vx)$, show that
  for the crew the explosion happened about 318 years before they passed you.
  What would an observer flying away from Betelgeuse at the same speed say?
\item \diff{1} A clock travels at $v = 0.99$ for one year of Earth time. How much
  time does it record?
\item \diff{2} \textbf{Accelerating twin.} Repeat the twins' trip with a more
  realistic flight plan: the traveller accelerates uniformly from rest to $0.8$
  over the first year of Earth time, cruises, and decelerates symmetrically at
  the star and on the way home. Set up the proper-time integral and evaluate it
  numerically.
\item \diff{2}\label{prob:gps-curves} \textbf{The curves of \cref{fig:clock-offsets}.} For a
  clock on a circular orbit of radius $r$, show that its daily gain over a clock
  on the ground is proportional to $1/R - 3/(2r)$. At what orbital radius does it
  vanish? Plot the gain from height, the loss from speed and the net result.
\item \diff{2}\label{prob:neutron-star} \textbf{Neutron star.} Show that a clock
  on the surface of a neutron star of mass $1.4\,\Msun$ and radius
  \SI{12}{\kilo\metre} runs at 0.81 of the rate of a distant clock. Use
  $G\Msun/c^{2} = \SI{1.477}{\kilo\metre}$.
\item \diff{2} \textbf{Polar coordinates.} For the flat two-dimensional metric
  $\dd s^{2} = \dd r^{2} + r^{2}\dd\phi^{2}$, compute all the Christoffel symbols
  from \cref{eq:christoffel}. Show that the geodesic equation has straight lines
  as solutions.
\item \diff{2} \textbf{Flowing coordinates, again.} Compute the Christoffel
  symbols of \cref{eq:flow-metric} and show that the Riemann tensor vanishes.
  Then replace $v_{0}$ by a function $v(y)$ and show that the Riemann tensor no
  longer vanishes.
\item \diff{2}\label{prob:geodesic} \textbf{Deriving the geodesic equation.}
  Write the proper time along a worldline as
  $\tau = \int \sqrt{-g_{\mu\nu}\dot x^{\mu}\dot x^{\nu}}\,\dd\lambda$, where the dot
  is $\dd/\dd\lambda$. Apply the Euler--Lagrange equations and choose
  $\lambda = \tau$ to obtain \cref{eq:geodesic}.
\item \diff{2}\label{prob:thrown-ball} \textbf{The thrown ball.} A ball is thrown
  straight up at speed $v_{0}$ and caught at the same height. Using
  \cref{eq:clock-rate} with $\Phi = gh/c^{2}$, show that its clock gains
  $v_{0}^{3}/(3gc^{2})$ over a clock held at the height of the hand. Then show that
  any other path $h(t)$ with the same endpoints gains less.
\item \diff{3}\label{prob:orbit-paradox} \textbf{The orbit paradox.} An
  astronaut in a circular orbit of radius $r$ passes a colleague standing on
  top of a tower of the same height, and they meet again one orbit later.
  Using \cref{eq:clock-rate}, show that the colleague on the tower, who is not
  in free fall, records more time, by the fraction $GM/(2rc^{2})$, and find the
  difference over one orbit at the height of the space station,
  $r = \SI{6790}{\kilo\metre}$. Then resolve the paradox. (a)~Show that two
  circular orbits of the same radius, tilted slightly with respect to each
  other, cross again after half an orbit, so that a full orbit is longer than
  the \enquote{short enough stretch} of the key idea of
  \topicref{sec:geodesics}. (b)~Describe a third worldline, also in free fall,
  that joins the same two events and records more time than either.
\item \diff{2} \textbf{The river.} Show that \cref{eq:river} turns into
  \cref{eq:schwarzschild} under the change of time coordinate
  $\dd T = \dd t + \frac{\sqrt{r_{\mathrm{s}}/r}}{1 - r_{\mathrm{s}}/r}\,\dd r$.
\item \diff{3} \textbf{Tides near a black hole.} Estimate the stretching
  acceleration across a \SI{2}{\metre} tall person falling feet first into a black
  hole of $10\,\Msun$ at the moment they cross the horizon, and again for a black
  hole of $10^{9}\,\Msun$. Which one could a traveller cross unharmed, and why
  does the answer depend on the mass the way it does?
\end{problems}

\begin{furtherreading}
\item[\textcite{Hartle2003}] A physics-first introduction to general relativity
  that builds intuition for metrics, clocks and geodesics before the formalism.
\item[\textcite{Schutz2009}] A careful, gentle development of spacetime
  diagrams, tensors and curvature, well suited to self-study.
\item[\textcite{Carroll2004}] A graduate-level introduction that covers the
  geometry of this chapter with full mathematical care.
\item[\textcite{Misner1973}] The encyclopedic classic, with deep geometric
  intuition; its sign conventions are the ones used in this book.
\item[\textcite{Wald1984}] General relativity in coordinate-free form.
\item[\textcite{Ashby2003}] Relativity in satellite navigation, including the
  Earth's rotation and shape.
\item[\textcite{Pound1960}] The first laboratory measurement of the
  gravitational redshift, in a tower at Harvard.
\item[\textcite{Chou2010}] Clocks that detect a change of height of
  \SI{33}{\centi\metre}.
\item[\textcite{Hamilton2008}] The river model of black holes, developed in
  full.
\item[\textcite{Michell1784}] Michell's dark stars, the first discussion of
  objects from which light cannot escape.
\end{furtherreading}
